\documentclass[letterpaper,11pt]{article}

\usepackage[utf8]{inputenc}
\usepackage{geometry}
\usepackage[hidelinks]{hyperref}
\usepackage{microtype}
\usepackage{helvet}
\renewcommand{\familydefault}{\sfdefault}

\geometry{
  left=0.9in,
  right=0.9in,
  top=0.8in,
  bottom=0.8in
}

\pagestyle{empty}
\linespread{1.08}
\setlength{\parindent}{0pt}
\setlength{\parskip}{10pt}

\begin{document}

%----------HEADER----------
\begin{center}
    {\LARGE \bfseries DARREL DANADYAKSA POLI} \\[3pt]
    \small Depok, Indonesia \quad $|$ \quad
    \href{mailto:darrel.danadyaksa19@gmail.com}{darrel.danadyaksa19@gmail.com} \quad $|$ \quad
    \href{https://www.linkedin.com/in/darrel-danadyaksa-poli/}{linkedin.com/in/darrel-danadyaksa-poli}
\end{center}
\vspace{2pt}
\hrule
\vspace{14pt}

%----------DATE----------
\today

\vspace{8pt}

Dear Hiring Manager,

I am writing to apply for an AI Research position. I am a final-year Computer Science student at the University of Indonesia with research experience at NVIDIA, where I investigated the theoretical and empirical properties of sequence architectures, and I am eager to bring that same rigor to your research team.

As an AI Research Intern at NVIDIA, I evaluated non-Transformer recurrent architectures (RWKV-v4 and RWKV-v7) against Transformers, LSTMs, and GRUs for formal regular language expressivity, and implemented a custom CUDA C++ WKV operator in PyTorch to accelerate recurrent state transition computations. To probe model decision boundaries rigorously, I designed synthetic formal-language benchmarks spanning parity counting, substring matching, and language unions, augmented with edit-distance hard-negative perturbations — work that required precise experimental design as much as implementation skill. Earlier, as part of the Indonesia AI x OMDENA Forest Fire project, I curated and normalized multi-source environmental datasets and fine-tuned predictive models for early wildfire risk detection.

This research instinct extends into my independent and competition work: I fine-tuned IndoBERT for Indonesian clickbait headline classification (Top 10 Finalist, Gemastik 2023), and built a multimodal architecture fusing BERT, Whisper, and VAD features via ViT-LSTM for emotion detection, achieving 99.71\% public leaderboard accuracy at Satria Data 2025. I am comfortable working across the research stack — from formal problem framing and benchmark design to CUDA-level implementation — in Python, PyTorch, TensorFlow, and Hugging Face.

I would welcome the opportunity to discuss how my background in architecture benchmarking, systems-level implementation, and applied model research could contribute to your research agenda. Thank you for your time and consideration.

\vspace{10pt}

Sincerely, \\
Darrel Danadyaksa Poli

\end{document}
